\documentclass[11pt]{article}
\usepackage{fontspec}
\usepackage[french,english]{babel}
\usepackage{amsmath}
\usepackage{amssymb}
\usepackage{amsthm}
\usepackage{mathrsfs}
\usepackage{tikz-cd}
\usepackage{enumitem}

\newcommand{\printedpage}[1]{\marginpar{\footnotesize p.~#1}}
\newcommand{\msp}[1]{\marginpar{\footnotesize [#1]}}
\newcommand{\Ch}{\mathrm{Ch}}
\newcommand{\Hot}{\mathrm{Hot}}
\newcommand{\Hom}{\mathrm{Hom}}

\begin{document}

\printedpage{[unnumbered]}
% head: (aucun titre courant imprimé)

\begin{center}
{\LARGE\textsf{THE DERIVATORS}}

\bigskip

\textbf{ALEXANDRE GROTHENDIECK}

\bigskip\bigskip

\textbf{Chapter IX}

\textbf{Return to the categories of fractions $\boldsymbol{\mathcal{M}W^{-1}}$,\\
comparison with Quillen}
\end{center}

\clearpage

\printedpage{[unnumbered]}
% head: (aucun titre courant imprimé)

This text was deciphered and transcribed in \LaTeXe{} by M. Künzer. It was edited by M. Künzer, J. Malgoire, and G. Maltsiniotis. The transcription is as faithful as possible to the manuscript. For the few obvious corrections, or rare comments by the editors, as well as for the original numbering of the manuscript pages, typewriter characters \texttt{[typewriter]} in square brackets are used. A question mark in square brackets means that one is not sure of the preceding word.

This edition is provisional. Remarks, comments and corrections are welcome. Send a message to:

\textbf{[e-mail address removed]}

\hfill G. Maltsiniotis\hspace*{2cm}

\clearpage

\printedpage{1}
% head: A. Grothendieck : Dérivateurs, Ch. IX, Édition des Universités de Montpellier II et Paris VII
% foot: Version 12 mars 2002

\msp{page 1}

\section*{1\quad Functoriality of $\Ch^{\Sigma}_{\infty}(a,b)$ in $a$, $b$ in $\mathcal{M}\Sigma^{-1}$}

We are given a category $\mathcal{M}$ with \texttt{[an]} admissible set of arrows $\Sigma \subseteq \mathrm{Fl}(\mathcal{M})$, whence (taking e.g. $\underline{t} = (t_n)$, with $t_n = (\underbrace{\Delta_+\Delta_-}_{\tau_0})^{2n}$, and $\underbrace{t_{n+1}}_{=\,\tau_0^{2n+2}} \longrightarrow \underbrace{t_n}_{=\,\tau_0^{2n}}$ defined as $p * \mathrm{id}_{\tau_0^{2n}} * p$, where $p : \tau_0 \longrightarrow 1$ the degeneracy.) We then have a theory of the $\underline{\Ch}_{\Sigma,\infty}(\mathcal{M},a,b)$, which I denote simply $\Ch_\infty(a,b)$ as an object of Hot. Pairings
\[
\Ch_\infty(a,b) \times \Ch_\infty(b,c) \longrightarrow \Ch_\infty(a,c) \qquad \text{in Hot}
\]
with (we shall admit it) \textit{associativity} in Hot. This therefore defines pairings on the $\pi_0(\Ch_\infty(a,b)) \overset{\text{def}}{=} \pi_0(a,b)$, which is none other than the composition of arrows in $\mathcal{M}\Sigma^{-1}$. On the other hand, we have for $a,b \in \mathrm{Ob}\,\mathcal{M}$
\[
\Hom_{\mathcal{M}}(a,b) \longrightarrow \pi_0(\Ch_\infty(a,b)) = \pi_0(a,b),
\]
compatible with compositions. But we can ignore it for the moment, and look at $\mathrm{Ob}\,\mathcal{M}$, or simply $\mathcal{M}_0$, with only the structure
\[
\left\{
\begin{array}{ccccccc}
\mathcal{M}_0 \times \mathcal{M}_0 & \longrightarrow & \Hot & \qquad & (a,b) & \longmapsto & \Ch_\infty(a,b) \\[1ex]
\mathcal{M}_0 \times \mathcal{M}_0 \times \mathcal{M}_0 & \longrightarrow & \mathrm{Fl}\,\Hot & \qquad & (a,b,c) & \longmapsto & (\Ch_\infty(a,b) \times \Ch_\infty(b,c) \underbrace{\longrightarrow}_{\substack{\text{arrow}\\ \text{in Hot}}} \Ch_\infty(a,c))
\end{array}
\right.
\]
with condition \texttt{[of]} associativity, and also a condition on units:
\[
\Ch_\infty(a,a) \text{ is equipped with a } 1_a \in \pi_0\Ch_\infty(a,a) = \pi_0(a,a),
\]
playing the role of a ``two-sided identity'', in the following sense. Let us note that the pairings
\[
\begin{array}{ccccc}
\Ch_\infty(a,b) & \times & \Ch_\infty(b,c) & \longrightarrow & \Ch_\infty(a,c) \\[1ex]
P & \times & Q & \longrightarrow & R
\end{array}
\]
\msp{page 2}
define
\[
\begin{array}{ccc}
P & \longrightarrow & \underline{\Hom}_{\Hot}(Q,R) \\[1ex]
Q & \longrightarrow & \underline{\Hom}_{\Hot}(P,R),
\end{array}
\]
whence on passing to the $\pi_0$
\[
\begin{array}{ccc}
\pi_0(P) & \longrightarrow & \Hom_{\Hot}(Q,R) \\[1ex]
\pi_0(Q) & \longrightarrow & \Hom_{\Hot}(P,R),
\end{array}
\]
\printedpage{2}
% head: A. Grothendieck : Dérivateurs, Ch. IX, Édition des Universités de Montpellier II et Paris VII
% foot: Version 12 mars 2002
in other words,
\[
\begin{array}{ccc}
\pi_0(a,b) & \longrightarrow & \Hom_{\Hot}(\Ch_\infty(b,c),\Ch_\infty(a,c)) \\[1ex]
\pi_0(b,c) & \longrightarrow & \Hom_{\Hot}(\Ch_\infty(a,b),\Ch_\infty(a,c)).
\end{array}
\]

We deduce from this in particular for $a,b \in \mathcal{M}_0$
\[
\left\{
\begin{array}{cccc}
1) & \pi_0(a,a) & \longrightarrow & \Hom_{\Hot}(\Ch_\infty(a,b),\Ch_\infty(a,b)) \\[1ex]
2) & \pi_0(a,a) & \longrightarrow & \Hom_{\Hot}(\Ch_\infty(b,a),\Ch_\infty(b,a)),
\end{array}
\right.
\tag*{$(*)$}
\]
and an object $\alpha_a$ of $\pi_0(a,a)$ is a two-sided unit if its two images are the identity homomorphisms in $\Ch_\infty(a,b)$ resp. $\Ch_\infty(b,a)$. Passing to the $\pi_0$ of these, this implies that $\alpha_a$ is a two-sided unit for the composition of the $\pi_0(x,y)$, relative to $b$. If we want this to be the case for \textit{every $b$,} this means (taking $b = a$) that $\alpha_a$ is a two-sided unit object for the structure of associative composition (hence of a monoid) in $\pi_0(a,a)$, hence that $\pi_0(a,a)$ is a monoid \textit{with unit,} having $\alpha_a$ as unit object. The element $\alpha_a$ is therefore uniquely determined. And we want $\alpha_a$ to exist
\msp{page 3}
for every $a \in \mathcal{M}_0$, so $\mathcal{M}_0$, equipped with the $\pi_0(a,b)$ and their composition, becomes a \textit{category,} say $\mathcal{M}_\Pi$; and its identity arrows must play the role of a two-sided unit also for the ``external Hom'' $\Ch_\infty(a,b)$, with values in Hot.

We can say that for fixed $a = a_0$, the $\Ch_\infty(a_0,b)$ ($b$ variable) form the values of a covariant \textit{functor}
\[
\mathcal{M}_\Pi \longrightarrow \Hot \qquad b \longmapsto \Ch_\infty(a_0,b),
\]
by associating to
\[
b \xrightarrow{\ u\ } b', \qquad u \in \pi_0(b,b'),
\]
the corresponding arrow
\[
\Ch_\infty(a,b) \longrightarrow \Ch_\infty(b') % [sic] \Ch_\infty(b') -> \Ch_\infty(a_0,b'), argument manquant (cf. la note des éditeurs qui suit)
\]
\texttt{[rather} $\Ch_\infty(a_0,b) \longrightarrow \Ch_\infty(a_0,b')$\texttt{]} $*_1$ [?]. That this is indeed a functorial law should follow formally from the associativity of the pairings on the $\Ch_\infty$. Likewise, we have a functorial law, for fixed $b$,
\[
\mathcal{M}^o_\Pi \longrightarrow \Hot, \qquad a \longmapsto \Ch_\infty(a,b),
\]
thanks to associativity. In fact, associativity will also imply that we have a \textit{bifunctor}
\[
\begin{array}{ccc}
\mathcal{M}^o_\Pi \times \mathcal{M}_\Pi & \longrightarrow & \Hot \\[1ex]
(a\,,\,b) & \longmapsto & \Ch_\infty(a,b),
\end{array}
\]
\printedpage{3}
% head: A. Grothendieck : Dérivateurs, Ch. IX, Édition des Universités de Montpellier II et Paris VII
% foot: Version 12 mars 2002
\msp{page 4}
for if
\[
a' \xrightarrow{\ u\ } a, \qquad b \xrightarrow{\ v\ } b',
\]
we have the diagram
\[
\begin{tikzcd}[column sep=large, row sep=large]
\Ch_\infty(a,b) \arrow[r, "\substack{\text{comp. with}\\ v}"] \arrow[d, "\text{comp. with } u"'] & \Ch_\infty(a,b') \arrow[d, "\text{comp. with } u"] \\
\Ch_\infty(a',b) \arrow[r, "\substack{\text{comp. with}\\ v}"] & \Ch_\infty(a',b'),
\end{tikzcd}
\]
which is commutative thanks to associativity for
\[
\Ch_\infty(a',a),\ \Ch_\infty(a,b),\ \Ch_\infty(b,b').
\]

\textbf{Corollary} of this bifunctorial dependence: \textit{If $u : a' \longrightarrow a$, $v : b \longrightarrow b'$ are in $\mathcal{M}_\Pi$, then the corresponding arrow}
\[
\Ch_\infty(a,b) \longrightarrow \Ch_\infty(a',b') \qquad \textit{in } \Hot
\]
\textit{is an} isomorphism.

The same arguments should work if instead of the $\Hom\,\mathrm{ext}$ with values in Hot, we took them in Ind\,Hot, or in Cat\,hot, or in Ind\,Cat\,hot -- I am thinking of course of the $\mathrm{CH}_\infty(a,b)$ (in the theory of localization) in Ind\,Cat\,hot, which can be regarded as coming from a bifunctor (where $\mathcal{M}_\Pi = \mathcal{M}\Sigma^{-1}$)
\[
\mathcal{M}_\Pi \times \mathcal{M}_\Pi \longrightarrow \mathrm{Ind}\,\mathrm{Cat}\,\mathrm{hot}.
\]

\msp{page 5}
Here is the definitive axiomatization (?). We replace Hot by a category

\begin{itemize}[label={$\mathcal{H}$,}]
\item category with product $*$ associative, commutative (?) (\footnote{We should not need commutativity.}), unital, and all the usual compatibilities as needed,
\end{itemize}

the unit object for $*$ will be denoted 1, and we shall write
\[
\left\{
\begin{array}{ccc}
\pi_0(\xi) & \overset{\text{def}}{=} & \Hom_{\mathcal{H}}(1,\xi) \\[1ex]
\pi_0 : \mathcal{H} & \longrightarrow & \mathrm{Ens}.
\end{array}
\right.
\]

If we have
\[
P * Q \xrightarrow{\ u\ } R,
\]

\end{document}
